% Généré par scripts/make_v11_prose.py, ne pas éditer à la main.

\subsection*{Période 2008-01 à 2024-01}

Aux États-Unis, sur les huit portefeuilles long short top~10, 6 affichent un TCAC positif. Le meilleur est le Hist Gradient Boosting, avec 7,4~\% par an (Sharpe 0,80, perte maximale 18,1~\%). Sur la même période, le S\&P 500 rend 7,6~\% (Sharpe 0,46) et le portefeuille équipondéré des titres de l'univers 10,8~\% (Sharpe 0,64) ; 3 portefeuille(s) long short sur huit font mieux que l'équipondéré en Sharpe.

Au Canada, sur les huit portefeuilles long short top~10, 2 affichent un TCAC positif. Le meilleur est le XGBoost, avec 5,6~\% par an (Sharpe 0,39, perte maximale 35,4~\%). Sur la même période, le composite TSX rend 2,6~\% (Sharpe 0,24) et le portefeuille équipondéré des titres de l'univers 11,5~\% (Sharpe 0,78) ; 0 portefeuille(s) long short sur huit font mieux que l'équipondéré en Sharpe.

Les $R^2$ hors échantillon moyens des régressions restent négatifs ou nuls (de $-$0,55 à 0,00 côté américain) : les prédictions expliquent moins bien les rendements qu'une moyenne simple, et la valeur des portefeuilles vient du classement, non du niveau prédit. Les modèles dont les prédictions sont identiques pour tous les titres à la plupart des dates (voir \texttt{results/v2/tables/prediction\_ties.csv}) doivent se lire comme des portefeuilles quasi fixes.

\subsection*{Période 2008-01 à 2012-01}

Aux États-Unis, sur les huit portefeuilles long short top~10, 7 affichent un TCAC positif. Le meilleur est le XGBoost, avec 10,8~\% par an (Sharpe 0,72, perte maximale 25,0~\%). Sur la même période, le S\&P 500 rend $-$3,8~\% (Sharpe 0,01) et le portefeuille équipondéré des titres de l'univers 4,4~\% (Sharpe 0,30) ; 7 portefeuille(s) long short sur huit font mieux que l'équipondéré en Sharpe.

Au Canada, sur les huit portefeuilles long short top~10, 2 affichent un TCAC positif. Le meilleur est le XGBoost de classification, avec 7,7~\% par an (Sharpe 0,53, perte maximale 13,4~\%). Sur la même période, le composite TSX rend $-$3,6~\% (Sharpe $-$0,01) et le portefeuille équipondéré des titres de l'univers 5,2~\% (Sharpe 0,35) ; 1 portefeuille(s) long short sur huit font mieux que l'équipondéré en Sharpe.

Les $R^2$ hors échantillon moyens des régressions restent négatifs ou nuls (de $-$0,46 à $-$0,00 côté américain) : les prédictions expliquent moins bien les rendements qu'une moyenne simple, et la valeur des portefeuilles vient du classement, non du niveau prédit. Les modèles dont les prédictions sont identiques pour tous les titres à la plupart des dates (voir \texttt{results/v2/tables/prediction\_ties.csv}) doivent se lire comme des portefeuilles quasi fixes.

\subsection*{Période 2012-01 à 2020-01}

Aux États-Unis, sur les huit portefeuilles long short top~10, 5 affichent un TCAC positif. Le meilleur est le XGBoost, avec 7,5~\% par an (Sharpe 0,90, perte maximale 15,1~\%). Sur la même période, le S\&P 500 rend 12,5~\% (Sharpe 0,99) et le portefeuille équipondéré des titres de l'univers 14,5~\% (Sharpe 1,17) ; 0 portefeuille(s) long short sur huit font mieux que l'équipondéré en Sharpe.

Au Canada, sur les huit portefeuilles long short top~10, 3 affichent un TCAC positif. Le meilleur est la régression logistique, avec 2,5~\% par an (Sharpe 0,24, perte maximale 24,7~\%). Sur la même période, le composite TSX rend 4,6~\% (Sharpe 0,47) et le portefeuille équipondéré des titres de l'univers 14,4~\% (Sharpe 1,42) ; 0 portefeuille(s) long short sur huit font mieux que l'équipondéré en Sharpe.

Les $R^2$ hors échantillon moyens des régressions restent négatifs ou nuls (de $-$0,16 à $-$0,04 côté américain) : les prédictions expliquent moins bien les rendements qu'une moyenne simple, et la valeur des portefeuilles vient du classement, non du niveau prédit. Les modèles dont les prédictions sont identiques pour tous les titres à la plupart des dates (voir \texttt{results/v2/tables/prediction\_ties.csv}) doivent se lire comme des portefeuilles quasi fixes.

\subsection*{Période 2020-01 à 2024-01}

Aux États-Unis, sur les huit portefeuilles long short top~10, 7 affichent un TCAC positif. Le meilleur est le XGBoost de classification, avec 6,5~\% par an (Sharpe 0,66, perte maximale 17,1~\%). Sur la même période, le S\&P 500 rend 10,3~\% (Sharpe 0,54) et le portefeuille équipondéré des titres de l'univers 9,4~\% (Sharpe 0,54) ; 4 portefeuille(s) long short sur huit font mieux que l'équipondéré en Sharpe.

Au Canada, sur les huit portefeuilles long short top~10, 3 affichent un TCAC positif. Le meilleur est le XGBoost, avec 6,9~\% par an (Sharpe 0,41, perte maximale 45,1~\%). Sur la même période, le composite TSX rend 5,3~\% (Sharpe 0,36) et le portefeuille équipondéré des titres de l'univers 11,6~\% (Sharpe 0,66) ; 0 portefeuille(s) long short sur huit font mieux que l'équipondéré en Sharpe.

Les $R^2$ hors échantillon moyens des régressions restent négatifs ou nuls (de $-$1,21 à $-$0,14 côté américain) : les prédictions expliquent moins bien les rendements qu'une moyenne simple, et la valeur des portefeuilles vient du classement, non du niveau prédit. Les modèles dont les prédictions sont identiques pour tous les titres à la plupart des dates (voir \texttt{results/v2/tables/prediction\_ties.csv}) doivent se lire comme des portefeuilles quasi fixes.
